\scanpage{41}
\mspage{- 27 -}
% archive sheet number: 40

\[
D\handwritten{u}(x,x_0,\handwritten{y}_0)=\handwritten{f}(x,\handwritten{u}(x,x_0,\handwritten{y}_0))\tag*{(20)}
\]
(D \uline{dériva tion par rapport à} x, $\handwritten{f}=(f_i)$), \uline{et la condition}
\[
\handwritten{u}(x_0,x_0,\handwritten{y}_0)=\handwritten{y}_0.\tag*{(21)}
\]

\uline{En outre}, \uline{il existe un voisinage} $V_1\handwritten{\subset}V$ \uline{de} $\handwritten{b}$ \uline{tel que}, \uline{pour tout} $\handwritten{y}_0\handwritten{\in}V_1$ \uline{et tout couple de points} $x,x_0$ \uline{de} U, \uline{l'équation} $\handwritten{y}_0=\handwritten{u}(x_0,x,\handwritten{y})$ \uline{ait dans} V \uline{une seule ra-cine} $\handwritten{y}=\handwritten{u}(x,x_0,\handwritten{y}_0)$.

Posons en effet $x-x_0=t$, $\handwritten{y}-\handwritten{y}_0=\handwritten{z}$ ; l'équation différentielle $\handwritten{y}'=\handwritten{f}(x,\handwritten{y})$ devient $\handwritten{z}'=\handwritten{f}(x_0+t,\handwritten{y}_0+\handwritten{z})$, et le second membre est holomorphe pour $(t,\handwritten{z})$ dans un voisinage de 0 dans $K^{p+1}$ et $(x_0,\handwritten{y}_0)$ dans un voisinage de $(a,\handwritten{b})$ ; la condition (21), pour une intégrale $\handwritten{v}$ de cette équation, s'écrit $\handwritten{v}(0;x_0,\handwritten{y}_0)=0$. On est donc dans les conditions d'application du th.2, d'où la première partie de la proposition. D'autre pa\handwritten{-}rt, si U a été pris assez petit, il existe un voisinage $V_1\handwritten{\subset}V$ de $\handwritten{b}$ tel que, pour tout $x_0\handwritten{\in}U$ et tout $\handwritten{y}_0\handwritten{\in}V_1$, l'image de U par l'application $x\handwritten{\to u}(x,x_0,\handwritten{y}_0)$ est contenue dans V car si M est le maximum de $\handwritten{\|f}(x,\handwritten{y})\handwritten{\|}$ dans $U\handwritten{\times}C$, on a en vertu du th. des accroissements finis, $\handwritten{\|u}(x,x_0,\handwritten{y}_0)-\handwritten{y}_0\handwritten{\|}=\handwritten{\|u}(x,x_0,\handwritten{y}_0)-\handwritten{u}(x_0,x_0,\handwritten{y}_0)\handwritten{\|\leq}M\handwritten{|}x-x_0\handwritten{|}$ quels que soient x et $x_0$ dans U. Si $\handwritten{y}=\handwritten{u}(x,x_0,\handwritten{y}_0)\handwritten{\in}V$, la fonction $s\handwritten{\to u}(s,x,\handwritten{y})$ est donc définie dans U et prend la valeur $\handwritten{y}$ au point x ; elle est donc égale à $\handwritten{u}(s,x_0,\handwritten{y}_0)$ en raison du th.2, et on a donc $\handwritten{u}(x_0,x,\handwritten{y})=\handwritten{u}(x_0,x_0,\handwritten{y}_0)=\handwritten{y}_0$. D'ailleurs, si $\handwritten{\unsure{w|W}\in}V$ est tel que $\handwritten{y}_0=\handwritten{u}(x_0,x,\handwritten{\unsure{w|W}})$, l'intégrale $s\handwritten{\to u}(s,x,\handwritten{\unsure{w|W}})$ qui est définie dans U, prend la valeur $\handwritten{y}_0$ au point $x_0$,donc est égale à $\handwritten{u}(s,x_0,\handwritten{y}_0)$ ; en particulier, elle prend la valeur $\handwritten{y}$ au point x, ce qui montre que $\handwritten{y}=\handwritten{\unsure{w|W}}$ et achève la démonstration.
